%% Shared by main.tex (inside the Discussion section, AAAI) and by
%% main_acl.tex (as the unnumbered \section*{Limitations} that *ACL requires).
We note the following limitations of this study.
\begin{enumerate}[leftmargin=*,itemsep=1pt,topsep=3pt]
  \item \textbf{Scale.} A 173-instance corpus is small, a consequence of prioritizing per-instance construction rigor.
  \item \textbf{Language balance.} The English-to-Chinese split is 53:120, which may favor models pre-trained more heavily in one of the two languages.
  \item \textbf{Dominating categories.} Entity scope (n = 94) and metric definition (n = 63) dominate, against only n = 9 for recognition policy, which suggests that current models are not able to generate instances of the rarer categories in the first place.
  \item \textbf{Generator and evaluator bias.} Under API budget constraints we use the GPT-5 family for construction and simulation, which may bias results toward that family.
  \item \textbf{Data contamination.} Instances are freshly constructed from FY2022-2025 facts, but we cannot rule out that recently released models were pre-trained on those filings.
\end{enumerate}
Appendix~\ref{app:limitations} expands each of these.
